\documentclass[main]{subfiles}


\title[Advanced Hyperparameter Optimization]{Advanced Hyperparameter Optimization}
\subtitle{Multi-Fidelity Optimization}
\author[Katharina Eggensperger]{Katharina Eggensperger}
\date{}

\titlebackground*{images/AutoML_Speeding_up_HPO.jpg}

%\institute{}
%\date{}

\begin{document}
	
	\maketitle
	

%----------------------------------------------------------------------
%----------------------------------------------------------------------
\begin{frame}[c]{Can't We Discard Bad Performing Configurations Early?}

$\leadsto$ Full evaluations yield the most accurate signal, but may waste resources in practice.

\vspace{0.5cm}

In practice, we often have access to cheap approximations or intermediate results of the target function, which can be useful for identifying promising regions and discarding poor-performing candidates early, e.g.,
\begin{itemize}
    \item subsampling training data
    \item partial training
    \item downscaled images
    \item smaller models
\end{itemize}

\vspace{0.5cm}
\textbf{$\leadsto$ How can we use this information to increase the \\ efficiency of Bayesian optimization?} \\

\end{frame}


%-----------------------------------------------------------------------
%----------------------------------------------------------------------
\begin{frame}[c]{A Motivating Example}

\begin{figure}
    \centering
    \includegraphics[width=0.8\textwidth]{05_Advanced_HPO/images/multi-fidelity_example.png}
    \caption{From: \liturl{Franceschi et al. "Hyperparameter Optimization in Machine Learning"}{https://arxiv.org/abs/2410.22854}}
\end{figure}

\end{frame}
%-----------------------------------------------------------------------
%----------------------------------------------------------------------
\begin{frame}[c]{Formal Definition}

We refer to an approximation of the true objective $\cost(\conf)$ as $\cost_b(\conf)$, where $b \in [b_{min}, b_{max}]$. Furthermore, we assume that $\cost_{b_{max}}(\conf) = \cost(\conf)$ and that $b$ correlates with evaluation costs (e.g. compute costs).

\vspace{1cm}
Our optimization objective is then to find:

\begin{equation*}
\conf^* \in \argmin_{\conf \in \pcs} \cost_{b_{max}}(\conf)
\end{equation*}

This requires allocating resources efficiently across configurations.
\end{frame}
%-----------------------------------------------------------------------
%----------------------------------------------------------------------
\begin{frame}[c]{Assumptions}

Multi-fidelity methods exploit three structural observations
\begin{enumerate}
    \item the \textbf{error shrinks smoothly} as the fidelity budget increases,
    \item \textbf{ranks are approximately consistent} across fidelities,
    \item and \textbf{higher fidelities provide a more accurate approximation} of the true objective value.
\end{enumerate}

\vspace{0.5cm}
$\leadsto$ In practice, these assumptions are not necessarily satisfied. HPO must be robust against that.

\end{frame}
%-----------------------------------------------------------------------
%----------------------------------------------------------------------
\begin{frame}{Successive Halving}

\textbf{Idea:} Explore configurations with low budgets and exploit promising configurations with high budgets.

\vspace{0.2cm}

\textbf{Method:} Sample $n$ random configurations $\confs_{cand}$: 
\begin{enumerate}
    \item evaluate $\confs_{cand}$ on budget $b$
    \item keep top $1/\eta$ fraction of $\confs_{cand}$
    \item increase budget $b \leftarrow \eta b$
    \item repeat until only one configuration remains or the maximum budget is reached.
\end{enumerate}

\vspace{0.2cm}

\textbf{Key parameter:} The reduction factor $\eta > 1$ controls how quickly we discard configurations and increase budget.
\begin{itemize}
    \item $\eta=2$ $\leadsto$ successively \emph{halves} the \#configurations; \emph{doubles} the budget.
    \item Larger $\eta$ $\leadsto$ aggressive pruning, greedy exploitation
    \item Smaller $\eta$ $\leadsto$ gentle pruning, guarded exploitation
\end{itemize}
\end{frame}
%-----------------------------------------------------------------------
%----------------------------------------------------------------------
\begin{frame}[c]{Successive Halving Limitations}

\textbf{Sensitivity to early decisions:} Pruning is based on low-fidelity evaluations. \\
$\leadsto$ Risk of discarding good configurations due to noise or a too low budget.

\vspace{0.3cm}

\textbf{No adaptive sampling:} Configurations are sampled randomly. \\
$\leadsto$ Does not learn from previous evaluations.

\vspace{0.3cm}

\textbf{Extensions:}
\begin{itemize}
    \item Successive halving with varying values for $b_{min}$ and $n$ $\leadsto$ Hyperband \liturl{Li et al. JMLR'18}{https://www.jmlr.org/papers/v18/16-558.html}
    \item Combine with Bayesian Optimization $\leadsto$ BOHB \liturl{Falkner et al. ICML'18}{https://proceedings.mlr.press/v80/falkner18a/falkner18a.pdf}
\end{itemize}

\end{frame}
%-----------------------------------------------------------------------
%----------------------------------------------------------------------
\begin{frame}{Hyperband}

\textbf{Motivation:} A small $b_{min}$ enables fast progress, if low-fidelity evaluations preserve ranking. The optimal value depends on the task and is typically unknown.

\vspace{0.3cm}

\begin{columns}[T,onlytextwidth]
  \column{0.48\textwidth}
  \textbf{Idea:} Run Successive Halving with different resource allocation schedules. In each \emph{bracket}, use the same total budget. \\
  $\leadsto$ Large $n$, small $b_{min}$ \hfill (broad exploration) \\
  $\leadsto$ Small $n$, large $b_{min}$ \hfill (deep evaluation) \\

  \vspace{0.3cm}
  \textbf{Intuition:} At least one bracket is close to optimal allocation, and the method is robust against budget constraints.

  \column{0.50\textwidth}
  \vspace{-0.5cm}
  \begin{figure}
      \centering
      \includegraphics[width=\linewidth]{05_Advanced_HPO/images/SH_pseudocode.png}
      \caption{From: \liturl{Falkner et al. JMLR'18}{proceedings.mlr.press/v80/falkner18a/falkner18a.pdf}}
  \end{figure}
\end{columns}

\end{frame}
%-----------------------------------------------------------------------
%----------------------------------------------------------------------
\begin{frame}[c]{Model-Based Multi-Fidelity Optimization}

\textbf{Motivation:} Replace random sampling with Bayesian Optimization

\vspace{0.3cm}

\textbf{Idea:} Use past observations to bias the selection of new configurations. Once there are sufficient observations per fidelity, replace random sampling with Bayesian Optimization\footnote{The original method uses a Tree-Parzen Estimator method for \\sampling-based Bayesian Optimization. See \liturl{BOHB}{https://proceedings.mlr.press/v80/falkner18a/falkner18a.pdf} and \liturl{TPE}{https://proceedings.neurips.cc/paper_files/paper/2011/hash/86e8f7ab32cfd12577bc2619bc635690-Abstract.html}.}.

\vspace{0.5cm}
\textbf{Intuition:} Focus search on promising regions.

\end{frame}
%-----------------------------------------------------------------------
%----------------------------------------------------------------------
\begin{frame}{Empirical Results}

\begin{figure}
    \centering
    \includegraphics[width=0.7\linewidth]{05_Advanced_HPO/images/BOHB.png}
    \caption{From \liturl{Falkner et al. JMLR'18}{https://proceedings.mlr.press/v80/falkner18a/falkner18a.pdf}}
\end{figure}
\end{frame}
%-----------------------------------------------------------------------
%----------------------------------------------------------------------
\begin{frame}[c]{Alternatives and Extensions}

Modern methods can further exploit \emph{parallelism} and \emph{structure}:

\begin{itemize}
    \item \textbf{Asynchronous} methods remove the synchronization barrier in Hyperband.
    \item \textbf{Learning Curve} methods predict performance of configurations from partial evaluations.
    \item \textbf{Scaling Laws} methods exploit predictable trends as a function of compute or data to extrapolate even more reliably from small budgets.
\end{itemize}

\end{frame}
%-----------------------------------------------------------------------
%----------------------------------------------------------------------
\begin{frame}{Conclusion and Key Takeaway}

Multi-fidelity methods improve the efficiency of hyperparameter optimization by allocating resources adaptively, e.g., by
\begin{itemize}
    \item evaluating many configurations with small budgets and progressively pruning (Successive Halving),
    \item hedging over multiple resource allocation strategies (Hyperband),
    \item combining adaptive resource allocation with model-based sampling (BOHB).
\end{itemize}

\vspace{0.5cm}

Performance depends on the quality of low-fidelity signals and the choice of resource allocation.
\end{frame}
%-----------------------------------------------------------------------
%----------------------------------------------------------------------

\end{document}